\section{Dispatch Optimization and Trading Simulation}\label{ch:methodology:dispatch-optimization-and-trading-simulation}
\subsection{Simulation Setup}
The simulation is executed on an hourly grid with $\Delta t = 1$ h. The battery state is represented by $SoC_t$ with lower and upper bounds $SoC_{\min}$ and $SoC_{\max}$, initial state $SoC_0$, and terminal target floor $SoC_T \ge SoC_{\text{target}}$. The converter power limit is $P_{\max}$ and charge/discharge efficiencies are $\eta_{\text{in}}$ and $\eta_{\text{out}}$.

\begin{equation}\label{eq:setup-soc-bounds}
    SoC_{\min} \le SoC_t \le SoC_{\max}
\end{equation}
\begin{equation}\label{eq:setup-soc-initial-terminal}
    SoC_0 = SoC_{\text{init}},
    \qquad
    SoC_T \ge SoC_{\text{target}}
\end{equation}
\begin{equation}\label{eq:setup-power-bounds}
    0 \le P_{\text{charge},t} \le P_{\max},
    \qquad
    0 \le P_{\text{discharge},t} \le P_{\max}
\end{equation}

Forecast inputs (ex-ante) are used for optimization; realized market data is only used in ex-post settlement.

\subsection{Rolling Dispatch Optimization}
At each decision time, a continuous linear optimization over horizon $H$ (default 48 h, smoke runs 24 h) is solved and only the first step is executed.
\begin{equation}\label{eq:rolling-dispatch-policy}
    \mathbf{u}^{*}_{t:t+H-1}
    =
    \arg\max_{\mathbf{u}} \hat{\Pi}_{t:t+H-1}(\mathbf{u}\mid\mathcal{F}_t),
    \qquad
    \mathbf{u}^{\mathrm{exec}}_{t}=\mathbf{u}^{*}_{t}
\end{equation}

Decision variables are continuous DA charge/discharge, directional reserve bids, reserve-bin allocations, auxiliary power and slack variables. The implementation uses an LP relaxation.

The objective combines expected DA/ID and aFRR value with soft-constraint penalties:
\begin{equation}\label{eq:dispatch-objective-soft}
    \max \hat{\Pi}^{\mathrm{LP}}
    =
    \sum_{t=1}^{T}\left(\hat{\pi}_{t}^{\mathrm{DA/ID}}+\hat{\pi}_{t}^{\mathrm{aFRR}}\right)
    - \lambda \sum_{t=1}^{T}\left(s_{t}^{+}+s_{t}^{-}\right)
\end{equation}

The ex-ante SoC transition used in optimization is:
\begin{equation}\label{eq:dispatch-soc-dynamics}
    SoC_{t+1}
    =
    SoC_t
    +
    \eta_{\text{in}}\,P_{\text{charge},t}\Delta t
    -
    \frac{P_{\text{discharge},t}\Delta t}{\eta_{\text{out}}}
    -
    E^{\text{aux}}_t
    -
    E^{\text{exp,act+}}_t
    +
    E^{\text{exp,act-}}_t
\end{equation}
where expected activation terms are represented through reserve-bin acceptance probabilities and activation-rate forecasts in the LP coefficients.

Converter coupling constraints for multi-market participation are:
\begin{equation}\label{eq:dispatch-power-coupling-neg}
    P_{\text{charge},t} + R^{-}_t \le P_{\max}
\end{equation}
\begin{equation}\label{eq:dispatch-power-coupling-pos}
    P_{\text{discharge},t} + R^{+}_t \le P_{\max}
\end{equation}
\begin{equation}\label{eq:dispatch-reserve-total}
    R^{+}_t + R^{-}_t \le R_{\max}
\end{equation}

Quantile-based reserve-feasibility constraints are implemented as soft chance constraints:
\begin{equation}\label{eq:dispatch-soft-p90-pos}
    SoC_t - \sum_{b\in\mathcal{B}}
    \left(
    \hat{p}^{+}_{\mathrm{acc},b,t}\,\hat{r}^{+}_{\tau,t}\,P^{+}_{b,t}\,\Delta t/\eta_{\text{out}}
    \right)
    + s_t^{+}
    \ge SoC_{\min}
\end{equation}
\begin{equation}\label{eq:dispatch-soft-p90-neg}
    SoC_t + \sum_{b\in\mathcal{B}}
    \left(
    \hat{p}^{-}_{\mathrm{acc},b,t}\,\hat{r}^{-}_{\tau,t}\,P^{-}_{b,t}\,\Delta t\,\eta_{\text{in}}
    \right)
    - s_t^{-}
    \le SoC_{\max}
\end{equation}

\subsection{Quantile Sweep}
The same rolling optimization model is re-run for different forecast quantile assumptions. Let $\tau$ denote the selected quantile. Quantile-indexed forecast inputs are:
\begin{equation}\label{eq:qs-da}
    \hat{\lambda}^{DA,\tau}_t
\end{equation}
\begin{equation}\label{eq:qs-afrr}
    \hat{C}^{aFRR,cap,\tau}_t,\;
    \hat{\lambda}^{aFRR,act,\tau}_t,\;
    \hat{r}^{aFRR,\tau}_t
\end{equation}
Each $\tau$ produces one strategy variant by solving Eq.~\ref{eq:rolling-dispatch-policy} with the corresponding quantile paths.

\subsection{Intraday Rebalancing and Recourse Logic}
Intraday (ID) recourse is modeled as a reactive one-hour-lag correction step to maintain physical feasibility and reduce non-delivery risk.
\begin{equation}\label{eq:id-lag}
    E^{ID}_{\mathrm{exec},t+1}
    =
    f\!\left(SoC^{\mathrm{post}}_t,\;R^{+}_{t+1},\;R^{-}_{t+1}\right)
\end{equation}

ID rescue power is limited by residual converter headroom after DA setpoints:
\begin{equation}\label{eq:id-headroom}
    P^{ID}_{\mathrm{ch},t}\le P_{\max}-P^{DA}_{\mathrm{ch},t},
    \qquad
    P^{ID}_{\mathrm{dis},t}\le P_{\max}-P^{DA}_{\mathrm{dis},t}
\end{equation}
This recourse layer is not used as an ex-ante arbitrage channel in the optimizer.

\subsection{Backtesting Procedure}
Backtesting is chronological and causal. At each step $t$: (i) use forecasts available at $t$, (ii) solve rolling optimization over $[t,t+H-1]$, (iii) execute first action, (iv) settle with realized outcomes, (v) update state, (vi) move to $t+1$. Formally:
\begin{equation}\label{eq:backtest-loop}
    \left(\mathcal{F}_t,SoC_t\right)
    \rightarrow
    \mathbf{u}^{*}_{t:t+H-1}
    \rightarrow
    \mathbf{u}^{\mathrm{exec}}_t
    \rightarrow
    \left(\text{realized settlement}_t, SoC_{t+1}\right)
\end{equation}

Previously closed DA commitments are fixed in subsequent re-optimizations (gate-closure lock-in), and model, quantile, and baseline variants are evaluated on the same realized timeline.
